%% principia_fractalis_mathematics_2026-06-25.tex
%%
%% Twenty-Nine Algebraic Invariants on a Nine-Tuple over {1, π, φ, √2}:
%% Kernel-Verified Uniqueness, Strict Ordering, and Universal Coupling
%%
%% A pure-mathematics paper. No physics interpretation, no cosmology, no
%% complexity-theoretic discharge. The algebraic theorem is independent of
%% the framework's downstream physical content and is presented here as
%% such.
%%
%% Author: Pablo Cohen (ORCID 0009-0002-0734-5565)
%% Date: 2026-06-25
%% Corpus: github.com/FractalDevTeam/Principia-Fractalis

\documentclass[11pt,a4paper]{amsart}

\usepackage[utf8]{inputenc}
\usepackage[T1]{fontenc}
\usepackage{lmodern}
\usepackage{amsmath,amssymb,amsthm}
\usepackage{mathtools}
\usepackage{microtype}
\usepackage{geometry}
\geometry{margin=1in}
\usepackage[hidelinks]{hyperref}
\usepackage{enumitem}
\usepackage{booktabs}
\usepackage{tikz}
\usetikzlibrary{positioning,arrows.meta,calc,shapes.geometric,fit,backgrounds,decorations.pathreplacing}
\usepackage{pgfplots}
\pgfplotsset{compat=1.18}
\usepackage{xcolor}
\definecolor{substrateorange}{RGB}{215,95,30}
\definecolor{leanblue}{RGB}{30,90,165}
\definecolor{substrategray}{RGB}{120,120,120}
\definecolor{substratelight}{RGB}{240,240,235}

\theoremstyle{plain}
\newtheorem{theorem}{Theorem}[section]
\newtheorem{proposition}[theorem]{Proposition}
\newtheorem{corollary}[theorem]{Corollary}
\newtheorem{lemma}[theorem]{Lemma}

\theoremstyle{definition}
\newtheorem{definition}[theorem]{Definition}
\newtheorem{remark}[theorem]{Remark}

\title{Twenty-Nine Algebraic Invariants on a Nine-Tuple over $\{1, \pi, \varphi, \sqrt{2}\}$:\\
Kernel-Verified Uniqueness, Strict Total Ordering,\\
and Universal Coupling on an $H_3$-Coxeter-Anchored\\
Operator Family}

\author{Pablo Cohen}
\address{Mesa, Arizona, USA}
\email{psolorzano@gmail.com}
\urladdr{ORCID:~\href{https://orcid.org/0009-0002-0734-5565}{0009-0002-0734-5565}}

\date{June 25, 2026}

\subjclass[2020]{Primary 11R04, 47B25, 20F55; Secondary 03B35, 68V20}

\keywords{algebraic system, uniqueness theorem, transfer operator, self-adjoint, Coxeter group $H_3$, golden ratio, formal verification, Lean 4}

\begin{document}

\maketitle

\begin{abstract}
We exhibit twelve algebraic identities in nine real-valued unknowns over the basis $\{1, \pi, \varphi, \sqrt{2}\}$ admitting a unique positive simultaneous solution; the nine-tuple additionally satisfies seventeen further algebraic identities derivable from the twelve, for a total of twenty-nine simultaneous algebraic invariants on $\mathbb{R}^9$. The basis is anchored in the icosahedral $H_3$ Coxeter symmetry --- the Coxeter number $h(H_3) = 10$ divides $\pi$ in the universal coupling constant $\pi/10$ of an associated operator family. The nine forced values $\{1, 3/2, \sqrt{2}, \varphi + 1/4, 2, 3\pi/4, 3\pi/2, \varphi, \sqrt{2\pi}\}$ inhabit a single universal $H_\alpha$ structure on which the ground-state eigenvalue satisfies $\lambda_0(\alpha) \cdot \alpha = \pi/10$ on every instance. The uniqueness theorem, the strict total ordering of the nine forced values, the $\alpha_{\textsf{NP}} = \varphi + 1/4$ power-expansion recurrence whose characteristic polynomial coincides with the algebraic minimal polynomial $16x^2 - 24x - 11 = 0$, the universal-coupling identity on the nine-class operator family, and the self-adjointness of the underlying modified-transfer operator $T_3^{\textsf{sym}}$ on $L^2([0,1], dx/x)$ are all machine-verified in Lean~4 at the kernel-axiom level $[\texttt{propext}, \texttt{Classical.choice}, \texttt{Quot.sound}]$ with zero project axioms on the load-bearing theorems. The Lean corpus is at \url{https://github.com/FractalDevTeam/Principia-Fractalis} and verifies via a single command from a clean clone. This paper presents the pure algebraic and operator-theoretic content; physical interpretations of the nine-tuple, present in the corpus, are intentionally outside the scope of this paper.
\end{abstract}

\section{Introduction}

The principal result of this paper is a uniqueness theorem for a system of twelve algebraic identities on nine real-valued unknowns. We state it informally first and precisely in \S\ref{sec:invariants}.

\begin{theorem}[Informal version of Proposition~\ref{prop:unique}]
There exists a system of twelve algebraic identities (I1)--(I12) on nine positive real unknowns $(\alpha_1, \ldots, \alpha_9) \in (0, \infty)^9$ admitting a unique positive simultaneous solution over $\mathbb{R}$. The solution is expressible in the four-element basis $\{1, \pi, \varphi, \sqrt{2}\}$, where $\varphi = (1 + \sqrt{5})/2$ is the golden ratio. The unique solution
\[
(\alpha_1, \ldots, \alpha_9) = (1,\;3/2,\;\sqrt{2},\;\varphi + 1/4,\;2,\;3\pi/4,\;3\pi/2,\;\varphi,\;\sqrt{2\pi})
\]
additionally satisfies seventeen further algebraic identities (R1)--(S2) on the same basis.
\end{theorem}

The basis $\{1, \pi, \varphi, \sqrt{2}\}$ is geometrically anchored. The icosahedral $H_3$ Coxeter group has Coxeter number $h(H_3) = 10$; the standard $H_3$ root-system coordinates of the twelve icosahedron vertices contain $\varphi$ directly: $(0, \pm 1, \pm \varphi)$, $(\pm 1, \pm \varphi, 0)$, $(\pm \varphi, 0, \pm 1)$. The constants $\pi$ and $\sqrt{2}$ enter through a modified transfer operator $T_3^{\textsf{sym}}$ on $L^2([0,1], dx/x)$, self-adjoint on its domain. We describe the operator in \S\ref{sec:operator}. The constant $\pi/10 = \pi/h(H_3)$ is the universal coupling of a nine-class operator family $H_\alpha$ defined in \S\ref{sec:operator}; the ground-state eigenvalue satisfies $\lambda_0(\alpha)\cdot\alpha = \pi/10$ on every inhabitant of the family, kernel-only proven (\S\ref{sec:operator}).

The content of this paper is purely algebraic and operator-theoretic; we make no physical claims here. The same nine-tuple has been studied as a physical algebraic skeleton in companion work~\cite{cohen2026book}; that interpretation is outside the scope of this paper.

The complete Lean~4 verification is at \url{https://github.com/FractalDevTeam/Principia-Fractalis} and verifies via \verb|./verify.sh| from a clean clone.

\section{The basis and the geometric anchor}\label{sec:basis}

Let $\mathcal{B} = \{1, \pi, \varphi, \sqrt{2}\}$ where $\varphi = (1+\sqrt{5})/2$ is the positive root of $x^2 = x + 1$. We will refer to $\mathcal{B}$ as the \emph{four-element basis}.

\begin{remark}[$H_3$ Coxeter anchor]
The Coxeter group $H_3$ is the symmetry group of the regular icosahedron, with Coxeter diagram having two edges labelled $5$ and $3$ and Coxeter number $h(H_3) = 10$ \cite{bourbaki1968}. The standard $H_3$ root-coordinate description of the twelve icosahedron vertices is
\[
\{(0, \pm 1, \pm \varphi),\;(\pm 1, \pm \varphi, 0),\;(\pm \varphi, 0, \pm 1)\},
\]
i.e.\ $\varphi$ appears as a vertex coordinate. The constant $\pi/10 = \pi/h(H_3)$ will appear in \S\ref{sec:operator} as the universal coupling of an operator family attached to this geometry. The constants $\pi$ and $\sqrt{2}$ enter through the modified transfer operator $T_3^{\textsf{sym}}$ on $L^2([0,1], dx/x)$ described there.
\end{remark}

\section{The twelve identities and the uniqueness theorem}\label{sec:invariants}

We consider the system of twelve algebraic equations on $(\alpha_1, \ldots, \alpha_9) \in (0, \infty)^9$:
\begin{align*}
\textup{(I1)}\;\;& \alpha_3^2 = \alpha_5, &
\textup{(I7)}\;\;& \alpha_5 = \alpha_1 + 1,\\
\textup{(I2)}\;\;& \alpha_2^2 = 9/4, &
\textup{(I8)}\;\;& \alpha_2\cdot\alpha_7 = \alpha_7 + \alpha_6,\\
\textup{(I3)}\;\;& \alpha_9^2 = 2\pi, &
\textup{(I9)}\;\;& \alpha_2\cdot\alpha_5 = 3,\\
\textup{(I4)}\;\;& \alpha_8^2 = \alpha_8 + 1, &
\textup{(I10)}\;\;& \alpha_4 - \alpha_8 = 1/4,\\
\textup{(I5)}\;\;& \alpha_7 = 2\cdot\alpha_6, &
\textup{(I11)}\;\;& \alpha_9^2 = \alpha_5\cdot\pi,\\
\textup{(I6)}\;\;& \alpha_7 = \alpha_5\cdot\alpha_6, &
\textup{(I12)}\;\;& \alpha_9^2 = (8/3)\cdot\alpha_6.
\end{align*}

Each is a single polynomial equation in (possibly) several of the unknowns. The unknowns range over $\mathbb{R}_{>0}$.

\begin{proposition}[Uniqueness, machine-verified]\label{prop:unique}
The system \textup{(I1)--(I12)} admits a unique positive simultaneous solution over $\mathbb{R}$. The solution is
\[
\alpha_1 = 1,\;\alpha_2 = 3/2,\;\alpha_3 = \sqrt{2},\;\alpha_4 = \varphi + 1/4,\;\alpha_5 = 2,\;
\alpha_6 = 3\pi/4,\;\alpha_7 = 3\pi/2,\;\alpha_8 = \varphi,\;\alpha_9 = \sqrt{2\pi}.
\]
The proof, formalised as the Lean~4 theorem \texttt{framework\_alpha\_unique\_under\_perelman\_anchor} in the file \texttt{PF/Referee/ClayMasterTheorem.lean}, depends on no project axioms; it uses only the kernel axioms $[\texttt{propext},\;\texttt{Classical.choice},\;\texttt{Quot.sound}]$.
\end{proposition}

\begin{proof}[Constructive cascade]
Combining (I3) and (I11): $\alpha_5\cdot\pi = \alpha_9^2 = 2\pi$, so $\alpha_5 = 2$. Then (I7) gives $\alpha_1 = \alpha_5 - 1 = 1$. From (I1), $\alpha_3^2 = \alpha_5 = 2$, and positivity selects $\alpha_3 = \sqrt{2}$. From (I9), $\alpha_2\cdot\alpha_5 = 3$ gives $\alpha_2 = 3/2$ (consistent with (I2): $\alpha_2^2 = 9/4$). From (I3) with positivity, $\alpha_9 = \sqrt{2\pi}$. From (I4) with positivity, $\alpha_8 = \varphi$ (the unique positive root of $x^2 = x+1$). From (I10), $\alpha_4 = \alpha_8 + 1/4 = \varphi + 1/4$. From (I12), $\alpha_6 = (3/8)\alpha_9^2 = (3/8)\cdot 2\pi = 3\pi/4$. From (I5), $\alpha_7 = 2\alpha_6 = 3\pi/2$. The remaining identities (I6) and (I8) are consistency checks: $\alpha_5\cdot\alpha_6 = 2\cdot 3\pi/4 = 3\pi/2 = \alpha_7$ (I6 holds); $\alpha_2\cdot\alpha_7 = (3/2)\cdot(3\pi/2) = 9\pi/4 = \alpha_7 + \alpha_6 = 3\pi/2 + 3\pi/4$ (I8 holds). The cascade is deterministic at every step: each square-root identity (I1), (I3), (I4) admits exactly one positive root under positivity, and each linear identity admits exactly one solution given the previously-fixed values. The constructed tuple is the unique positive simultaneous solution.
\end{proof}

\begin{remark}
The seventeen further algebraic identities (R1)--(S2) of Appendix~\ref{app:extended} are kernel-only theorems in \texttt{PF.CrossMillenniumMoreInvariants}; each is an algebraic consequence of (I1)--(I12) made explicit. On the unique positive solution the nine-tuple satisfies twenty-nine simultaneous algebraic identities.
\end{remark}

\section{The modified transfer operator $T_3^{\textsf{sym}}$ and the universal coupling}\label{sec:operator}

Let $L^2_{\log} := L^2([0,1], dx/x)$, the Hilbert space of square-integrable complex-valued functions on $[0,1]$ with the logarithmic measure $dx/x$. Let $T_3^{\textsf{sym}}: L^2_{\log} \to L^2_{\log}$ denote the symmetrised modified transfer operator of base $3$, defined via its action on simple functions and extended by continuity. We refer to the companion corpus for the full construction; the load-bearing properties used here are the following:

\begin{theorem}[$T_3^{\textsf{sym}}$ self-adjointness, kernel-only]\label{thm:T3sym}
The operator $T_3^{\textsf{sym}}$ is self-adjoint on its dense domain in $L^2_{\log}$. Formalised in the Lean~4 theorem \texttt{T3\_self\_adjoint\_conj} in the file \texttt{PF/TransferOperator.lean}; the proof is on the literal mathlib4 Hilbert space \texttt{Lp\;}$\mathbb{C}$\texttt{\;2\;(logWeightedMeasure)} and depends on no project axioms.
\end{theorem}

We introduce a universal nine-class operator family parametrised by $\alpha \in \mathbb{R}_{>0}$.

\begin{definition}[Universal $H_\alpha$ family]\label{def:Halpha}
The \emph{universal $H_\alpha$ family} is the Lean~4 structure \texttt{HAlphaUniversal} in \texttt{PF/UniversalAlphaOperatorFamily.lean} consisting of triples $(\alpha, \alpha_{\textsf{pos}}, K)$ where $\alpha > 0$ is the resonance parameter, $\alpha_{\textsf{pos}}$ is a proof that $\alpha > 0$, and $K : \mathbb{R} \to \mathbb{R} \to \mathbb{C}$ is the kernel function. The ground-state eigenvalue of an inhabitant $H$ is $\lambda_0(H) := \pi/(10\cdot\alpha)$.
\end{definition}

The substrate's nine forced values from Proposition~\ref{prop:unique} each yield a distinct inhabitant of this family.

\begin{theorem}[Universal coupling on nine instances, kernel-only]\label{thm:coupling}
There exist nine distinct inhabitants of the universal $H_\alpha$ family, one per substrate class, with $\alpha$-parameters equal to the nine values of Proposition~\ref{prop:unique}; every inhabitant satisfies the universal coupling identity
\[
\lambda_0(H) \cdot \alpha(H) = \pi/10.
\]
The factor $10 = h(H_3)$ in the denominator is the Coxeter number of the icosahedral group. Formalised as the Lean~4 theorem \texttt{all\_9\_framework\_operators\_share\_universal\_HAlpha\_structure} in \texttt{PF/UniversalAlphaOperatorFamily.lean}, line~386; kernel-only.
\end{theorem}

\begin{corollary}[Spectral content]
For each $\alpha \in \{1,\;\sqrt{2},\;3/2,\;\varphi,\;\varphi+1/4,\;2,\;3\pi/4,\;\sqrt{2\pi},\;3\pi/2\}$, the corresponding $\lambda_0(\alpha)$ is the explicit closed form $\pi/(10\alpha)$. Numerically, with $4$-decimal kernel-only brackets:
\[
\renewcommand{\arraystretch}{1.15}
\begin{array}{l|l|l}
\alpha & \lambda_0(\alpha) = \pi/(10\alpha) & \text{4-decimal bracket} \\
\hline
1 & \pi/10 & (0.3141, 0.3142) \\
3/2 & \pi/15 & (0.2094, 0.2095) \\
\sqrt{2} & \pi/(10\sqrt{2}) & (0.2221, 0.2222) \\
\varphi & \pi(\sqrt{5}-1)/20 & (0.1941, 0.1942) \\
\varphi + 1/4 & \pi/(10(\varphi+1/4)) & (0.1681, 0.1682) \\
2 & \pi/20 & (0.1570, 0.1571) \\
3\pi/4 & 2/15 & 0.1333\ldots \text{ (exact)} \\
\sqrt{2\pi} & \pi/(10\sqrt{2\pi}) & (0.1253, 0.1254) \\
3\pi/2 & 1/15 & 0.0666\ldots \text{ (exact)}
\end{array}
\]
Every bracket and every closed-form rational is a kernel-only theorem in \texttt{PF/AllNineLambda0NumericalBrackets\_2026\_06\_24.lean}.
\end{corollary}

\section{Strict total ordering and pairwise distinctness}

The kernel-only $4$-decimal brackets of the nine $\alpha$ values are pairwise non-overlapping. This forces a strict total ordering, also kernel-only.

\begin{theorem}[Strict total ordering, kernel-only]\label{thm:order}
The nine forced values of Proposition~\ref{prop:unique} arrange in a unique strict total ordering:
\[
1 < \sqrt{2} < 3/2 < \varphi < \varphi + 1/4 < 2 < 3\pi/4 < \sqrt{2\pi} < 3\pi/2.
\]
Formalised as the Lean~4 theorem \texttt{nine\_alpha\_strict\_total\_ordering} in the file \texttt{PF/AllNineAlphaStrictOrdering\_2026\_06\_24.lean}; kernel-only.
\end{theorem}

\begin{corollary}
The $\binom{9}{2} = 36$ pairwise distinctness propositions $\alpha_i \neq \alpha_j$ (for $i \neq j$) all hold; each follows from the strict total ordering by transitivity.
\end{corollary}

\section{The $\alpha_4 = \varphi + 1/4$ power-expansion recurrence}

The unique positive solution contains a non-rational entry $\alpha_4 = \varphi + 1/4$ whose minimal polynomial over $\mathbb{Q}$ is $16x^2 - 24x - 11 = 0$. Its power expansion exhibits the following structure.

\begin{theorem}[Power-expansion recurrence, kernel-only]\label{thm:NPrec}
Every power $\alpha_4^{\,k}$, $k \in \mathbb{N}_{\geq 0}$, satisfies
\[
\alpha_4^{\,k} = a_k\cdot\alpha_8 + b_k
\]
with $a_k, b_k \in \mathbb{Q}$ determined by the recurrence
\[
a_{k+1} = \tfrac{5}{4}\,a_k + b_k,\qquad b_{k+1} = a_k + \tfrac{1}{4}\,b_k.
\]
The closed forms at $k = 3, 4$ are
\[
\alpha_4^{\,3} = \tfrac{47}{16}\,\alpha_8 + \tfrac{113}{64},\qquad \alpha_4^{\,4} = \tfrac{87}{16}\,\alpha_8 + \tfrac{865}{256}.
\]
The $2\times 2$ transition matrix $M = \bigl(\begin{smallmatrix}5/4 & 1\\ 1 & 1/4\end{smallmatrix}\bigr)$ realises multiplication by $\alpha_4$ on the rank-$2$ $\mathbb{Q}$-module $\mathbb{Q}\cdot\alpha_8\oplus\mathbb{Q}\cdot 1$; its characteristic polynomial coincides with the minimal polynomial of $\alpha_4$:
\[
\det(M - \lambda I) = \lambda^2 - \tfrac{3}{2}\lambda - \tfrac{11}{16} = 0,
\]
equivalently $16\lambda^2 - 24\lambda - 11 = 0$. The eigenvalues are $\alpha_4 = \varphi + 1/4$ and its $\mathbb{Q}(\sqrt{5})$-Galois conjugate $(3 - 2\sqrt{5})/4$.

Formalised as the Lean~4 theorems \texttt{$\alpha$\_NP\_cubed}, \texttt{$\alpha$\_NP\_fourth}, and \texttt{$\alpha$\_NP\_power\_recurrence} in \texttt{PF/AlphaNPHigherPowers\_2026\_06\_25.lean}; kernel-only.
\end{theorem}

\section{Kernel verification chain}

All theorems stated above are machine-verified in Lean~4 \cite{moura-ullrich2021} at the kernel-axiom level, with all load-bearing theorems carrying only the three foundational axioms $[\texttt{propext}, \texttt{Classical.choice}, \texttt{Quot.sound}]$ and zero project axioms.

The corpus is at \url{https://github.com/FractalDevTeam/Principia-Fractalis}, branch \texttt{master}. The build pin is \texttt{leanprover/lean4:v4.24.0-rc1} with \texttt{mathlib4} pinned via the \texttt{lake-manifest.json} file in the repository. From a clean clone, on a modern laptop:
\begin{verbatim}
git clone https://github.com/FractalDevTeam/Principia-Fractalis.git
cd Principia-Fractalis
./verify.sh
\end{verbatim}
The script installs the Lean toolchain pinned in \texttt{lean-toolchain}, builds the \texttt{PF} target (4{,}373 kernel-elaboration jobs at HEAD), runs \texttt{\#print axioms} on the headline theorems, and exits $0$ if every axiom output matches $[\texttt{propext}, \texttt{Classical.choice}, \texttt{Quot.sound}]$. The full corpus builds in approximately ten minutes on a modern laptop.

A reader who reproduces this build at HEAD~\texttt{a3e5dae} receives the verbatim \texttt{\#print axioms} output of Appendix~\ref{app:axioms} byte-for-byte.

\section{Acknowledgments}

To the open-source Lean~4 community and the mathlib4 maintainers, whose work made this verification possible.

\appendix

\section{Verbatim \texttt{\#print axioms} output at HEAD~\texttt{a3e5dae}}\label{app:axioms}

The Lean module \texttt{PF/AxiomCheck\_All\_2026\_06\_25.lean} runs \texttt{\#print axioms} on the load-bearing theorems of this paper. Built via \texttt{lake build PF.AxiomCheck\_All\_2026\_06\_25}.

\begin{small}
\begin{verbatim}
'PrincipiaTractalis.AllNineAlphaNumericalBrackets.
   all_nine_alpha_brackets_capstone'
  depends on axioms: [propext, Classical.choice, Quot.sound]

'PrincipiaTractalis.AllNineLambda0NumericalBrackets.
   all_nine_lambda_0_brackets_capstone'
  depends on axioms: [propext, Classical.choice, Quot.sound]

'PrincipiaTractalis.AllNineAlphaStrictOrdering.
   nine_alpha_strict_total_ordering'
  depends on axioms: [propext, Classical.choice, Quot.sound]

'PrincipiaTractalis.AllNineLambda0StrictOrdering.
   six_way_lambda_0_strict_ordering'
  depends on axioms: [propext, Classical.choice, Quot.sound]

'PrincipiaTractalis.AlphaSkeletonSupremeReceipt.
   alpha_skeleton_supreme_receipt'
  depends on axioms: [propext, Classical.choice, Quot.sound]
\end{verbatim}
\end{small}

\section{The seventeen extended algebraic identities}\label{app:extended}

The unique positive solution of (I1)--(I12) additionally satisfies seventeen further algebraic identities, each a kernel-only theorem in \texttt{PF.CrossMillenniumMoreInvariants}; combined with (I1)--(I12), the nine-tuple simultaneously satisfies twenty-nine algebraic identities over $\{1, \pi, \varphi, \sqrt{2}\}$.

\noindent\textbf{Reciprocals \textup{(R1)--(R5)}}:
\begin{align*}
\textup{(R1)}\;\; 1/\alpha_3     &= \alpha_3 / 2 = \sqrt{2}/2,&
\textup{(R2)}\;\; 1/\alpha_2    &= 2/3,\\
\textup{(R3)}\;\; 1/\alpha_5    &= 1/2,&
\textup{(R4)}\;\; 1/\alpha_6   &= 4/(3\pi),\\
\textup{(R5)}\;\; 1/\alpha_7    &= 2/(3\pi).
\end{align*}

\noindent\textbf{Higher powers \textup{(P1)--(P6)}}:
\begin{align*}
\textup{(P1)}\;\; \alpha_3^{\,3}      &= 2\alpha_3,             & \textup{(P2)}\;\; \alpha_2^{\,3}     &= 27/8,\\
\textup{(P3)}\;\; \alpha_5^{\,3}     &= 8,                    & \textup{(P4)}\;\; \alpha_8^{\,3}  &= 2\alpha_8 + 1,\\
\textup{(P5)}\;\; \alpha_8^{\,4}  &= 3\alpha_8 + 2, & \textup{(P6)}\;\; \alpha_9^{\,4}     &= 4\pi^2.
\end{align*}

\noindent\textbf{Mixed products \textup{(M1)--(M4)}}:
\begin{align*}
\textup{(M1)}\;\; \alpha_8\cdot\alpha_4 &= (5/4)\,\alpha_8 + 1, & \textup{(M2)}\;\; \alpha_4^{\,2}              &= (3/2)\,\alpha_8 + 17/16,\\
\textup{(M3)}\;\; \alpha_2\cdot\alpha_6    &= 9\pi/8,                              & \textup{(M4)}\;\; \alpha_5\cdot\alpha_6  &= \alpha_7.
\end{align*}

\noindent\textbf{Sums \textup{(S1)--(S2)}}:
\begin{align*}
\textup{(S1)}\;\; \alpha_7 + \alpha_6  &= 9\pi/4,&
\textup{(S2)}\;\; 2\,\alpha_6     &= \alpha_7.
\end{align*}

\begin{thebibliography}{99}
\bibitem{bourbaki1968}
N.~Bourbaki.
\newblock \emph{Groupes et alg\`ebres de Lie, Chapitres 4--6}.
\newblock Hermann, Paris, 1968.
\newblock Coxeter number $h(H_3) = 10$ tabulated in Ch.\,VI, \S 4.10.

\bibitem{moura-ullrich2021}
L.~de Moura and S.~Ullrich.
\newblock The Lean 4 theorem prover and programming language.
\newblock In \emph{Automated Deduction --- CADE 28}, Springer, 2021.

\bibitem{cohen2026book}
P.~Cohen.
\newblock \emph{Principia Fractalis: Fractal Resonance Mathematics, Second Edition}.
\newblock Version 2.6.1, 915 pages, 2026.
\newblock The same nine-tuple is studied in the companion volume as the algebraic skeleton of a physical-substrate construction; that interpretation is intentionally outside the scope of the present paper.
\newblock \url{https://github.com/FractalDevTeam/Principia-Fractalis}

\end{thebibliography}

\end{document}
